\documentclass[12pt,a4paper]{article}

\usepackage{fontspec}
\usepackage{polyglossia}
\setdefaultlanguage{russian}
\setmainfont{Times New Roman}
\setsansfont{Arial}
\setmonofont{Consolas}

\usepackage[a4paper,margin=22mm,headheight=16pt]{geometry}
\usepackage{microtype}
\usepackage{amsmath,amssymb,mathtools,bm}
\usepackage{booktabs,array,tabularx,longtable}
\usepackage{enumitem}
\usepackage{xcolor}
\usepackage{fancyhdr}
\usepackage{hyperref}

\definecolor{navy}{HTML}{17365D}
\definecolor{blue}{HTML}{2F75B5}
\definecolor{lightblue}{HTML}{EAF3F8}
\definecolor{green}{HTML}{548235}
\definecolor{lightgreen}{HTML}{EEF5E9}
\definecolor{orange}{HTML}{C55A11}
\definecolor{lightorange}{HTML}{FFF2E8}
\definecolor{red}{HTML}{A61C1C}
\definecolor{lightred}{HTML}{FCECEC}
\definecolor{lightgray}{HTML}{F3F4F6}

\hypersetup{colorlinks=true,linkcolor=navy,urlcolor=blue,citecolor=green,pdftitle={Active Block Wishart TPP: алгоритм с нуля},pdfauthor={Учебный конспект по реализации Wishart TPP}}

\pagestyle{fancy}
\fancyhf{}
\fancyhead[L]{\small Active Block Wishart TPP}
\fancyhead[R]{\small Учебный конспект}
\fancyfoot[C]{\thepage}
\renewcommand{\headrulewidth}{0.4pt}

\setlength{\parindent}{0pt}
\setlength{\parskip}{0.65em}
\setlist{topsep=0.35em,itemsep=0.25em,parsep=0.1em,leftmargin=1.7em}
\allowdisplaybreaks
\renewcommand{\arraystretch}{1.2}

\DeclareMathOperator{\tr}{tr}
\DeclareMathOperator{\diag}{diag}
\DeclareMathOperator{\softmax}{softmax}
\DeclareMathOperator{\Cat}{Cat}
\DeclareMathOperator{\Wishart}{Wishart}
\DeclareMathOperator{\KL}{KL}
\newcommand{\E}{\mathbb{E}}
\newcommand{\R}{\mathbb{R}}
\newcommand{\History}{\mathcal{H}}

\newenvironment{intuitionbox}{\par\medskip\noindent{\color{blue}\bfseries Интуиция. }\ignorespaces}{\par\medskip}
\newenvironment{examplebox}{\par\medskip\noindent{\color{orange}\bfseries Численный пример. }\ignorespaces}{\par\medskip}
\newenvironment{warningbox}{\par\medskip\noindent{\color{red}\bfseries Важно не перепутать. }\ignorespaces}{\par\medskip}
\newenvironment{summarybox}{\par\medskip\noindent{\color{navy}\bfseries Короткий итог. }\ignorespaces}{\par\medskip}

\newcommand{\term}[1]{\textbf{#1}}
\newcommand{\code}[1]{\texttt{#1}}

\begin{document}

\begin{titlepage}
\centering
\vspace*{20mm}
{\sffamily\bfseries\color{navy}\fontsize{28}{34}\selectfont Active Block Wishart TPP\par}
\vspace{5mm}
{\sffamily\Large Алгоритм, принцип работы и математика с нуля\par}
\vspace{14mm}
\fcolorbox{blue}{lightblue}{\begin{minipage}{0.82\textwidth}
Этот конспект последовательно строит модель от обычного пуассоновского процесса до смеси нейросетевых временных точечных процессов с матричным случайным эффектом. Предварительные знания ограничиваются школьной алгеброй; производные, интегралы, вероятности и матрицы вводятся по мере необходимости.
\end{minipage}}
\vfill
{\large Учебная версия, согласованная с текущей реализацией алгоритма\par}
\vspace{5mm}
{\large \today\par}
\end{titlepage}

\setcounter{tocdepth}{1}
\tableofcontents
\newpage

\section*{Как читать этот документ}
\addcontentsline{toc}{section}{Как читать этот документ}

Модель состоит из нескольких слоёв, и почти вся сложность возникает, когда их пытаются понять одновременно. Поэтому мы будем двигаться снизу вверх: сначала разберём отдельное событие и интенсивность, затем обычную смесь кластеров, после этого добавим индивидуальную скрытую матрицу и только в конце соберём полный алгоритм обучения.

В каждом разделе есть три уровня объяснения: словесная интуиция, формула и связь с программным кодом. Формулы не нужно запоминать. Важно понимать, какую задачу решает каждый объект и какие данные на него влияют.

\begin{summarybox}
Вся модель умещается в одной фразе: нейросеть предлагает для каждого кластера базовую интенсивность, скрытая матрица индивидуально преобразует её для конкретной траектории, а likelihood определяет, какой кластер и какое преобразование лучше объясняют наблюдаемые события.
\end{summarybox}

\section{С чего начинается задача: последовательность событий}

\subsection{События, время и метки}

Пусть мы наблюдаем некоторую систему во времени. В ней происходят события разных типов. Например, в медицинской истории типами могут быть диагноз, анализ и назначение лекарства; в журнале действий пользователя --- просмотр, клик и покупка; в сетевом журнале --- разные категории запросов.

Одну наблюдаемую последовательность, или \term{траекторию}, запишем как
\begin{equation}
s_m=\bigl\{(t_{m1},c_{m1}),\ldots,(t_{mN_m},c_{mN_m})\bigr\},
\end{equation}
где \(m\) --- номер траектории, \(t_{mi}\) --- время \(i\)-го события, \(c_{mi}\in\{1,\ldots,C\}\) --- его тип, а \(N_m\) --- число событий.

События упорядочены по времени: \(0<t_{m1}<\cdots<t_{mN_m}\leq T_m\). Величина \(T_m\) называется горизонтом наблюдения. Историю до момента \(t\) обозначим через \(\History_t\): это все события, которые уже успели произойти к этому моменту.

\begin{examplebox}
Если \(C=3\), последовательность \(\{(0.8,2),(1.1,1),(2.7,2)\}\) означает: в момент \(0.8\) произошло событие второго типа, в момент \(1.1\) --- первого, а в момент \(2.7\) --- снова второго. До времени \(t=2\) история содержит первые два события.
\end{examplebox}

\subsection{Что такое интенсивность}

Для каждого типа \(c\) модель задаёт положительную функцию \(\lambda_c(t\mid\History_t)\), которая называется \term{условной интенсивностью}. Она показывает мгновенную склонность системы породить событие типа \(c\) в момент \(t\), если известна вся доступная история.

Для очень маленького промежутка длины \(\Delta t\) выполняется приближение
\begin{equation}
\Pr\{\text{событие типа }c\text{ на }[t,t+\Delta t)\mid\History_t\}\approx \lambda_c(t\mid\History_t)\Delta t.
\end{equation}

Интенсивность не является вероятностью: она может быть больше единицы и измеряется в «событиях на единицу времени». Вероятность получается только после умножения на короткий интервал времени.

Полная интенсивность равна сумме по типам:
\begin{equation}
\lambda^{\mathrm{total}}(t\mid\History_t)=\sum_{c=1}^{C}\lambda_c(t\mid\History_t).
\end{equation}

Она управляет общим темпом событий. Отношение \(\lambda_c/\lambda^{\mathrm{total}}\) описывает, какая доля мгновенной активности приходится на тип \(c\).

\begin{intuitionbox}
Представьте несколько лампочек, по одной на каждый тип события. Яркость лампочки \(c\) равна \(\lambda_c(t)\). Суммарная яркость говорит, насколько скоро ожидается любое событие, а относительные яркости определяют наиболее вероятный тип.
\end{intuitionbox}

\subsection{Простейший случай: процесс Пуассона}

Если \(C=1\) и \(\lambda(t)=r\) постоянно, мы получаем однородный пуассоновский процесс. За интервал длины \(T\) он в среднем порождает \(rT\) событий. Например, при \(r=2\) и \(T=5\) ожидается \(10\) событий, хотя конкретная случайная траектория может содержать \(7\), \(10\) или \(14\) событий.

Нейросетевой временной точечный процесс обобщает эту идею: интенсивность меняется со временем и зависит от предыдущих событий. Backbone-модель --- например COTIC, THP, NHP или RMTPP --- вычисляет эту сложную функцию по данным.

\section{Как likelihood оценивает интенсивность}

\subsection{Два слагаемых с противоположными ролями}

Для траектории \(s_m\) логарифм правдоподобия многомерного временного точечного процесса равен
\begin{equation}
\ell_m=\sum_{i=1}^{N_m}\log \lambda_{c_{mi}}(t_{mi}\mid\History_{t_{mi}})-\int_0^{T_m}\sum_{c=1}^{C}\lambda_c(t\mid\History_t)\,dt.
\label{eq:tpp-loglik}
\end{equation}

Первая часть награждает модель за высокую интенсивность наблюдавшегося типа ровно в моменты событий. Вторая часть штрафует её за высокую интенсивность на всём интервале: нельзя безнаказанно объявить, что события чрезвычайно вероятны всегда.

\begin{intuitionbox}
Likelihood требует честного прогноза. Модель получает очки, если заранее подсветила фактические события, но платит за каждую область времени, где обещала много событий, а данных для этого обещания не было.
\end{intuitionbox}

\subsection{Почему возникает интеграл}

Разобьём время на очень маленькие интервалы. На интервале без события вероятность приблизительно равна \(1-\lambda^{\mathrm{total}}(t)\Delta t\). Если перемножить эти вероятности для всех пустых интервалов и перейти к пределу \(\Delta t\to0\), получится экспонента \(\exp\{-\int\lambda^{\mathrm{total}}(t)dt\}\). После логарифмирования экспонента превращается во второе слагаемое формулы~\eqref{eq:tpp-loglik}.

\subsection{Как интеграл считается на компьютере}

Для сложной нейросетевой интенсивности первообразной обычно нет. Поэтому интеграл оценивают численно. В текущем алгоритме доступны детерминированная квадратура Гаусса и Лежандра и метод Монте-Карло.

Для Монте-Карло на интервале \([a,b]\) генерируются независимые числа \(u_r\sim\mathrm{Uniform}(0,1)\), затем точки \(t_r=a+(b-a)u_r\), и используется оценка
\begin{equation}
\int_a^b \lambda^{\mathrm{total}}(t)\,dt\approx \frac{b-a}{R}\sum_{r=1}^{R}\lambda^{\mathrm{total}}(t_r).
\label{eq:mc-integral}
\end{equation}

Оценка~\eqref{eq:mc-integral} несмещённая: в среднем по случайным точкам она равна истинному интегралу. Во время обучения точки можно пересэмплировать, а на validation и test полезно фиксировать их, чтобы сравнение моделей не загрязнялось новым шумом Монте-Карло.

\section{Обычная смесь временных точечных процессов}

\subsection{Зачем нужны кластеры}

Предположим, что траектории принадлежат к \(K\) качественно различным группам. Например, одна группа пользователей действует редко, но целенаправленно, а другая --- часто и хаотично. Для траектории \(m\) вводится скрытая категориальная переменная
\begin{equation}
z_m\sim\Cat(\pi_1,\ldots,\pi_K),\qquad \pi_k\geq0,\qquad \sum_{k=1}^{K}\pi_k=1.
\end{equation}

Число \(z_m=k\) означает принадлежность к кластеру \(k\), а \(\pi_k\) --- априорную долю этого кластера в популяции.

Нейросетевой backbone выдаёт не одну, а \(K\) базовых интенсивностей:
\begin{equation}
h_{\theta,k}(t\mid\History_t)=\bigl(h_{\theta,k1}(t\mid\History_t),\ldots,h_{\theta,kC}(t\mid\History_t)\bigr)\in\R_+^C.
\end{equation}

Параметры нейросети обозначены через \(\theta\). В чистой смеси при \(z_m=k\) интенсивность просто равна \(h_{\theta,k}\).

\subsection{Как выбирается кластер}

Пусть \(\ell_{mk}\) --- логарифм likelihood траектории \(m\) под интенсивностью компоненты \(k\). Тогда posterior-вероятность кластера равна
\begin{equation}
\gamma_{mk}=\Pr(z_m=k\mid s_m)=\softmax_k\bigl(\log\pi_k+\ell_{mk}\bigr).
\end{equation}

Softmax превращает произвольные оценки в положительные числа с суммой один. Компонента получает большую ответственность, если она была вероятна заранее и хорошо объяснила события.

\begin{examplebox}
Пусть два кластера одинаково вероятны, а их log-likelihood равны \(-10\) и \(-12\). Отношение posterior-весов равно \(e^{-10}/e^{-12}=e^2\approx7.39\). Поэтому первый кластер получает вероятность около \(0.881\), а второй --- около \(0.119\).
\end{examplebox}

\subsection{Ограничение чистой смеси}

В чистой смеси все траектории одного кластера используют одну и ту же функцию \(h_{\theta,k}\). Их различает только обычный случайный шум процесса. Однако реальные траектории одного класса могут иметь устойчивые индивидуальные особенности: один субъект в среднем активнее другого, у них различаются пропорции типов событий, и эти различия сохраняются на всём горизонте.

Нейросеть может попытаться восстановить индивидуальность из истории, но тогда ей приходится хранить её неявно в конечном скрытом состоянии. При коротком контексте, разрывах последовательности или новом субъекте это неудобно. Нам нужна отдельная переменная, которая явно представляет постоянный эффект траектории.

\section{Индивидуальный случайный эффект}

\subsection{Три уровня неоднородности}

Active Block Wishart TPP разделяет вариативность на три уровня:
\begin{enumerate}
\item \(z_m\) описывает дискретный тип траектории, то есть различие между кластерами.
\item \(U_m\) описывает устойчивое индивидуальное отклонение внутри выбранного кластера.
\item Сам TPP создаёт случайные события даже при фиксированных \(z_m\) и \(U_m\).
\end{enumerate}

Таким образом, один кластер больше не обязан быть одной фиксированной интенсивностью. Он становится распределением над индивидуальными интенсивностями.

\subsection{Почему эффект является матрицей}

Скалярный random effect мог бы только одновременно ускорять или замедлять все типы событий. Диагональная матрица могла бы независимо менять их масштабы. Полная симметричная матрица дополнительно способна смешивать направления: индивидуальное усиление одного типа может быть связано с изменением другого.

В модели используется матрица \(U_m\in\R^{C\times C}\). Она действует только на \(C\) типов событий активного кластера, поэтому её размер не растёт с числом кластеров \(K\). Отсюда название \term{active block}: преобразуется активный \(C\)-мерный блок, выбранный переменной \(z_m\).

\section{Минимум линейной алгебры}

\subsection{Вектор и матрица}

Вектор \(x=(x_1,\ldots,x_C)^\top\) --- столбец чисел. Матрица \(A\in\R^{C\times C}\) преобразует его в новый вектор \(Ax\). Каждый выход является взвешенной суммой входов:
\begin{equation}
(Ax)_i=\sum_{j=1}^{C}A_{ij}x_j.
\end{equation}

Если \(A=I_C\), где \(I_C\) --- единичная матрица, то \(Ax=x\). Диагональная матрица только масштабирует координаты, а внедиагональные элементы смешивают их.

\subsection{Положительно определённая матрица}

Матрица \(U\) называется симметричной положительно определённой, или SPD, если \(U=U^\top\) и \(x^\top Ux>0\) для любого ненулевого \(x\). У SPD-матрицы положительные собственные значения, существует обратная матрица и однозначный главный симметричный квадратный корень \(U^{1/2}\), удовлетворяющий
\begin{equation}
U^{1/2}U^{1/2}=U.
\end{equation}

Если \(U=V\diag(d_1,\ldots,d_C)V^\top\), где столбцы \(V\) --- ортонормированные собственные векторы и \(d_i>0\), то
\begin{equation}
U^{1/2}=V\diag(\sqrt{d_1},\ldots,\sqrt{d_C})V^\top.
\end{equation}

След матрицы --- сумма диагональных элементов: \(\tr U=\sum_i U_{ii}\). Для SPD-матрицы он также равен сумме собственных значений.

\begin{intuitionbox}
SPD-матрицу удобно представлять как преобразование эллипса: собственные векторы задают главные направления, а собственные значения --- силу растяжения вдоль них. Внедиагональные элементы поворачивают главные направления относительно исходных типов событий.
\end{intuitionbox}

\section{Wishart-распределение}

\subsection{Распределение над матрицами}

Обычное нормальное распределение задаёт случайное число или вектор. Распределение Уишарта задаёт случайную SPD-матрицу. В нашей модели
\begin{equation}
U_m\mid z_m=k\sim\Wishart_C\!\left(\nu,\frac{\Omega_k}{\nu}\right),
\label{eq:wishart-prior}
\end{equation}
причём параметризация выбрана так, что
\begin{equation}
\E[U_m\mid z_m=k]=\Omega_k.
\end{equation}

Матрица \(\Omega_k\) описывает среднюю популяционную геометрию индивидуальных преобразований внутри кластера \(k\). Число \(\nu\) называется числом степеней свободы: при прочих равных большее \(\nu\) сильнее концентрирует \(U_m\) около \(\Omega_k\).

Чтобы общий масштаб нельзя было произвольно переносить между базовой интенсивностью и \(\Omega_k\), вводится ограничение
\begin{equation}
\tr\Omega_k=C.
\label{eq:trace-constraint}
\end{equation}

Ограничение~\eqref{eq:trace-constraint} фиксирует средний масштаб собственных значений, но оставляет модели возможность обучать их относительные величины и направления.

\begin{warningbox}
\(\Omega_k\) --- не скрытая матрица конкретной траектории. \(\Omega_k\) общая для популяции кластера, а \(U_m\) индивидуальна. У двух субъектов одного кластера один и тот же prior \(p(U\mid k)\), но разные случайные реализации \(U_m\).
\end{warningbox}

\section{Active-block преобразование интенсивности}

\subsection{Wishart-преобразование и convex-интерполяция}

Для матрицы \(U\) определим преобразование базовой интенсивности
\begin{equation}
T_U(h)=\left[U^{1/2}\sqrt h\right]^{\circ2}.
\label{eq:operator}
\end{equation}

Параметр \(\alpha\) обучается вместе с моделью. Он отвечает за то, насколько сильно индивидуальная матрица влияет на интенсивность.

При \(z_m=k\) итоговая интенсивность равна
\begin{equation}
\boxed{\lambda_m(t\mid z_m=k,U_m)=(1-\alpha)h_{\theta,k}(t\mid\History_t)+\alpha T_{U_m}\!\left(h_{\theta,k}(t\mid\History_t)\right)+\varepsilon,\quad\varepsilon=10^{-6}.}
\label{eq:active-rate}
\end{equation}

Корень и квадрат внутри \(T_U\) берутся покоординатно, а \(U_m^{1/2}\) --- матричный корень. Обозначение \(v^{\circ2}\) означает \((v_1^2,\ldots,v_C^2)\). Интерполяция между базовой и преобразованной ветвями выполняется уже в пространстве интенсивностей. Постоянный floor \(\varepsilon\) одинаков в событийном члене и компенсаторе.

\subsection{Почему сначала корень, а потом квадрат}

Базовая интенсивность \(h\) положительна. В Wishart-ветви мы переходим к амплитуде \(x=\sqrt h\), линейно преобразуем её матрицей \(U^{1/2}\), а затем возводим координаты в квадрат. Эта ветвь неотрицательна, после чего она convex-комбинируется с \(h\).

Это позволяет полной матрице смешивать типы событий, не нарушая главное требование TPP: \(\lambda_c(t)\geq0\).

\subsection{Точное вложение чистой смеси}

Если \(\alpha=0\), то
\begin{equation}
\lambda_m(t\mid z_m=k,U_m)=h_{\theta,k}(t)+\varepsilon.
\end{equation}

Матрица \(U_m\) полностью исчезает из likelihood. Следовательно, Pure Mixture является точным граничным случаем, а не приблизительной абляцией. Если random effect не нужен данным, модель может сама уменьшить \(\alpha\) до нуля.

При \(\alpha=1\) используется полное преобразование \(U^{1/2}\). Промежуточные значения дают плавную интерполяцию.

\subsection{Пошаговый численный пример}

\begin{examplebox}
Пусть есть два типа событий, а backbone выдал базовую интенсивность \(h=(4,1)^\top\). Тогда амплитуда равна \(\sqrt h=(2,1)^\top\). Возьмём \(U=\diag(2.25,0.25)\), поэтому \(U^{1/2}=\diag(1.5,0.5)\), и положим \(\alpha=0.4\). Wishart-ветвь равна
\begin{equation*}
T_U(h)=\left[\diag(1.5,0.5)(2,1)^\top\right]^{\circ2}=(9,0.25)^\top.
\end{equation*}
Итоговая интенсивность без учёта микроскопического floor равна \(0.6(4,1)^\top+0.4(9,0.25)^\top=(6,0.7)^\top\). Базовый общий темп был \(5\), а индивидуальный стал \(6.7\). Одновременно доля первого типа выросла с \(0.8\) примерно до \(0.896\). Значит, одна матрица изменила и темп, и состав меток.
\end{examplebox}

Если добавить внедиагональные элементы, каждая новая амплитуда будет зависеть сразу от нескольких базовых типов. Это и есть матричное перепутывание.

\section{Полная генеративная история}

Генеративная модель описывает мысленный процесс, которым данные могли быть созданы. Для каждой новой траектории \(m\) выполняются три шага:
\begin{enumerate}
\item Выбрать кластер \(z_m\sim\Cat(\pi)\).
\item Сэмплировать индивидуальную матрицу \(U_m\mid z_m=k\sim\Wishart_C(\nu,\Omega_k/\nu)\).
\item Сгенерировать события TPP с интенсивностью из формулы~\eqref{eq:active-rate}.
\end{enumerate}

\begin{center}
\renewcommand{\arraystretch}{1.4}
\begin{tabular}{c}
\fcolorbox{navy}{lightblue}{веса \(\pi\)} \(\longrightarrow\) \fcolorbox{green}{lightgreen}{кластер \(z_m\)} \(\longrightarrow\) \fcolorbox{green}{lightgreen}{эффект \(U_m\)} \(\longrightarrow\) \fcolorbox{navy}{lightblue}{интенсивность \(\lambda_m\)} \(\longrightarrow\) \fcolorbox{green}{lightgreen}{события \(s_m\)} \\
\(\Omega_{1:K}\nearrow\qquad h_{\theta,1:K}\nearrow\)
\end{tabular}
\end{center}

Стрелка от \(U_m\) к \(\lambda_m\) означает генеративное влияние. После наблюдения \(s_m\) направление рассуждения обращается: мы используем события, чтобы вычислить posterior по \(z_m\) и \(U_m\). Это не новая генеративная стрелка, а статистический вывод скрытых причин по наблюдаемому следствию.

\section{Likelihood компоненты с random effect}

При фиксированных \(z_m=k\) и \(U_m=U\) логарифм likelihood равен
\begin{equation}
\ell_{mk}(U)=\sum_{i=1}^{N_m}\log\lambda_{mk,c_{mi}}(t_{mi};U)-\int_0^{T_m}\sum_{c=1}^{C}\lambda_{mkc}(t;U)\,dt.
\label{eq:component-likelihood}
\end{equation}

В отличие от contrast-softmax конструкций, полная интенсивность \(\sum_c\lambda_{mkc}(t)\) зависит и от кластера, и от \(U\). Поэтому разница между траекториями с малым и большим числом событий непосредственно входит в likelihood через интегральный штраф и значения интенсивности в моменты событий.

Если бы индивидуальная матрица была известна, кластерные ответственности вычислялись бы обычным mixture likelihood. Но \(U_m\) скрыта, поэтому для нового субъекта нужно интегрировать её:
\begin{equation}
p(s_m\mid z_m=k)=\int p(s_m\mid z_m=k,U,\theta)\,p(U\mid k,\Omega_k)\,dU.
\label{eq:marginal-component}
\end{equation}

Интеграл~\eqref{eq:marginal-component} многомерный и содержит нейросеть и матричный корень, поэтому аналитического решения обычно нет.

\section{Вариационный вывод: приближаем posterior}

\subsection{Что такое posterior}

Prior \(p(U\mid k)\) описывает индивидуальную матрицу до наблюдения событий. Posterior \(p(U\mid s_m,k)\) учитывает конкретную траекторию. Чем информативнее события, тем сильнее posterior может отличаться от популяционного prior.

Точный posterior недоступен. Мы заменяем его удобным распределением того же семейства:
\begin{equation}
q_{mk}(U)=\Wishart_C\!\left(\kappa_{mk},\frac{Q_{mk}}{\kappa_{mk}}\right),\qquad \E_{q_{mk}}[U]=Q_{mk}.
\label{eq:variational-family}
\end{equation}

Для каждой пары «траектория--компонента» существуют свои параметры \(Q_{mk}\) и \(\kappa_{mk}\). Это важно: пока кластер неизвестен, одна и та же траектория рассматривается под каждой возможной гипотезой \(k\).

Матрица \(Q_{mk}\) параметризуется как \(L_{mk}L_{mk}^\top\), где \(L_{mk}\) нижнетреугольная с положительной диагональю. Такая Cholesky-параметризация гарантирует положительную определённость во время градиентной оптимизации.

\subsection{Free energy}

Качество вариационного posterior для пары \((m,k)\) измеряется free energy:
\begin{equation}
\mathcal F_{mk}=-\E_{q_{mk}}[\ell_{mk}(U)]+\KL\!\left(q_{mk}(U)\,\|\,p(U\mid k)\right).
\label{eq:free-energy}
\end{equation}

Первое слагаемое требует хорошо объяснять наблюдаемые события. Второе удерживает индивидуальный posterior около популяционного prior, если данных недостаточно для уверенного отклонения.

\begin{intuitionbox}
KL-дивергенцию можно представить как резинку между \(q_{mk}\) и \(p_k\). Likelihood тянет локальную матрицу к объяснению конкретной траектории, а резинка не позволяет изобрести произвольный \(U_m\), который идеально подгонит данные, но не относится к популяции.
\end{intuitionbox}

Ожидание \(\E_q[\ell]\) оценивается сэмплированием матриц \(U\sim q\). Для дифференцирования используется Bartlett-reparameterization: случайный шум отделяется от обучаемых параметров, и градиент может пройти через сэмпл.

\subsection{Responsibilities с неизвестной матрицей}

Free energy заменяет недоступный логарифм маргинального likelihood. Ответственности равны
\begin{equation}
\gamma_{mk}=\softmax_k\bigl(\log\pi_k-\mathcal F_{mk}\bigr).
\label{eq:responsibilities}
\end{equation}

Маленькая free energy означает, что компонента может объяснить траекторию правдоподобной индивидуальной матрицей. Большая free energy возникает, если нужен плохой likelihood или слишком далёкий от population law posterior.

\begin{examplebox}
Пусть три компоненты имеют одинаковые \(\pi_k\), а free energy равны \(10\), \(12\) и \(15\). Тогда веса пропорциональны \(e^{-10}\), \(e^{-12}\) и \(e^{-15}\), а ответственности приблизительно равны \(0.876\), \(0.119\) и \(0.006\). Первая компонента почти полностью объясняет траекторию.
\end{examplebox}

\subsection{ELBO как общая цель}

Полная вариационная нижняя граница имеет вид
\begin{equation}
\mathcal L=\sum_{m=1}^{M}\sum_{k=1}^{K}\gamma_{mk}\left[\E_{q_{mk}}\ell_{mk}(U)-\KL(q_{mk}\|p_k)+\log\pi_k-\log\gamma_{mk}\right].
\label{eq:elbo}
\end{equation}

Внутри скобок находятся четыре понятные части: качество TPP likelihood, штраф за неправдоподобный individual effect, prior кластера и энтропия responsibilities. Обучение стремится увеличить \(\mathcal L\), то есть уменьшить отрицательный ELBO.

\section{Как обучается популяционная матрица \texorpdfstring{\(\Omega_k\)}{Omega}}

\subsection{Достаточная статистика}

После локального inference у каждой пары \((m,k)\) есть средняя матрица \(Q_{mk}\) и ответственность \(\gamma_{mk}\). Их взвешенное среднее внутри компоненты равно
\begin{equation}
S_k=\frac{\sum_m\gamma_{mk}Q_{mk}}{\sum_m\gamma_{mk}}.
\label{eq:sufficient-statistic}
\end{equation}

Если бы ограничения на след не было, естественным обновлением было бы \(\Omega_k=S_k\). Но при \(\tr\Omega_k=C\) требуется решить constrained-задачу.

\subsection{Точный M-шаг}

Зависящая от \(\Omega_k\) часть Wishart KL приводит к задаче
\begin{equation}
\Omega_k^{\star}=\arg\min_{\Omega\succ0,\,\tr\Omega=C}\left\{\log|\Omega|+\tr(\Omega^{-1}S_k)\right\}.
\label{eq:omega-mstep}
\end{equation}

Оптимум имеет те же собственные векторы, что и \(S_k\). Если \(s_i\) --- собственные значения \(S_k\), а \(\omega_i\) --- искомые собственные значения \(\Omega_k\), условие оптимальности даёт
\begin{equation}
s_i=\omega_i+\xi\omega_i^2,\qquad \sum_{i=1}^{C}\omega_i=C,
\label{eq:eigen-kkt}
\end{equation}
где \(\xi\) --- один множитель Лагранжа для ограничения на след.

Алгоритм решает скалярные уравнения~\eqref{eq:eigen-kkt}, перебирает допустимые положительные ветви при необходимости, выбирает решение с минимальным objective и собирает матрицу обратно из собственных векторов. После этого используется небольшой eigenvalue floor для численной устойчивости и повторная нормировка следа.

\begin{warningbox}
Эвристика \(\Omega=CS/\tr S\) лишь масштабирует \(S\), но в общем случае не решает задачу~\eqref{eq:omega-mstep}. Trace constraint меняет собственные значения нелинейно. Такая проекция может исказить раннюю геометрию и заставить нейросеть компенсировать ошибку.
\end{warningbox}

\subsection{Почему обновление называется exact}

Слово exact относится не ко всей модели, а к условному M-шагу по \(\Omega\): при фиксированных \(q\), \(\gamma\), \(\theta\), \(\alpha\) и \(\pi\) constrained-минимум находится спектральным решателем, а не несколькими приближёнными градиентными шагами. Вариационный posterior и интегралы likelihood всё равно остаются приближёнными.

\section{Обновления \texorpdfstring{\(\pi\)}{pi} и \texorpdfstring{\(\alpha\)}{alpha}}

\subsection{Веса смеси}

При фиксированных responsibilities обновление весов имеет простой вид
\begin{equation}
\pi_k\leftarrow\frac{1}{M}\sum_{m=1}^{M}\gamma_{mk}.
\end{equation}

То есть новая априорная доля кластера равна средней posterior-ответственности за него.

\subsection{Сила random effect}

Для \(\alpha\) при фиксированных posterior-сэмплах минимизируется
\begin{equation}
J(\alpha)=-\sum_{m,k}\gamma_{mk}\frac{1}{R}\sum_{r=1}^{R}\ell_{mk}\bigl(U_{mk}^{(r)};\alpha\bigr),\qquad 0\leq\alpha\leq1.
\label{eq:alpha-objective}
\end{equation}

Поскольку \(\alpha\) --- один скаляр, удобнее bounded one-dimensional search, чем Adam. Реализация использует golden-section search и обязательно сравнивает внутренний кандидат с \(\alpha=0\), \(\alpha=1\) и предыдущим значением.

Если данные не поддерживают случайный эффект, граничная точка \(\alpha=0\) должна иметь не худшее значение objective. Если индивидуальная неоднородность полезна, оптимум остаётся положительным.

\begin{warningbox}
Маленькое \(\alpha\) не всегда означает маленькую функциональную вариативность: данные видят композицию \(\alpha(U^{1/2}-I)\), а масштаб отклонения \(U\) частично компенсирует \(\alpha\). Поэтому интерпретировать следует итоговые интенсивности и predictive quality, а не один скаляр отдельно.
\end{warningbox}

\section{Обновление нейросетевого backbone}

При фиксированных \(q\), \(\gamma\), \(\Omega\), \(\alpha\) и \(\pi\) нейросеть обучается максимизировать взвешенный ожидаемый log-likelihood:
\begin{equation}
\max_{\theta}\sum_{m,k}\gamma_{mk}\E_{q_{mk}}\left[\ell_{mk}(U;\theta)\right].
\label{eq:neural-mstep}
\end{equation}

Responsibilities и сэмплы \(U\) на этом шаге считаются константами; градиент идёт только в банк базовых интенсивностей. Backbone не обязан быть COTIC. Обёртке достаточно интерфейса, который возвращает положительный банк \(h_{\theta,1:K}\) и значения для event term и compensator.

После изменения \(\theta\) прежний локальный posterior становится немного устаревшим. Поэтому алгоритм чередует несколько разных шагов, а не обучает все параметры одним неструктурированным градиентом.

\section{Полный цикл обучения}

\subsection{Почему используется чередование}

Параметры тесно связаны: \(q(U)\) зависит от текущей нейросети и \(\Omega\), \(\Omega\) оценивается по \(q(U)\), а нейросеть обучается при текущих \(U\). Чередующийся variational EM разделяет ответственности и делает каждый подшаг понятным и проверяемым.

Один базовый цикл Active Wishart выполняется так:
\begin{enumerate}
\item Заморозить backbone и построить cache его базовых интенсивностей для train-траекторий.
\item Выполнить локальный E-шаг: оптимизировать \(q_{mk}(U)\) градиентом по free energy.
\item Вычислить \(\gamma_{mk}\) по формуле~\eqref{eq:responsibilities}; в нескольких первых циклах можно временно балансировать компоненты, чтобы разрушить симметрию смеси.
\item Выполнить exact constrained M-шаг по \(\Omega_k\).
\item Обновить веса \(\pi_k\) средними responsibilities.
\item Выполнить ограниченный одномерный поиск по \(\alpha\).
\item Сделать несколько neural M-шагов по формуле~\eqref{eq:neural-mstep} с persistent Adam optimizer.
\item Оценить validation prior-predictive NLL и сохранить лучший checkpoint.
\end{enumerate}

\medskip
\noindent{\color{navy}\bfseries Псевдокод полного цикла}
\begin{center}
\begin{tabularx}{\textwidth}{@{}>{\bfseries}p{8mm}X@{}}
0 & Инициализировать \(\theta,\Omega_{1:K},\pi,\alpha\).\\
1 & Для каждого outer cycle построить trace cache при текущем backbone.\\
2 & Несколько шагов минимизировать \(\mathcal F_{mk}\) по параметрам \(q_{mk}\).\\
3 & Посчитать responsibilities \(\gamma_{mk}\).\\
4 & Из \(Q_{mk}\) и \(\gamma_{mk}\) собрать \(S_k\), затем точно обновить \(\Omega_k\).\\
5 & Обновить \(\pi_k\) и найти лучшее \(\alpha\in[0,1]\).\\
6 & Несколько minibatch-шагов обновлять только \(\theta\).\\
7 & Посчитать validation prior-predictive NLL; при улучшении сохранить состояние.\\
8 & После последнего цикла вернуть лучший validation-checkpoint, а не последний.\\
\end{tabularx}
\end{center}

В базовой конфигурации для all-12 DAN Pure и Wishart обучаются как отдельные эксперименты от общего короткого checkpoint. Wishart использует \(75\) outer cycles по \(8\) neural updates, то есть \(600\) neural minibatch-шагов; Pure получает тот же neural budget. Локальный inference использует \(10\) Adam-шагов и \(2\) Wishart-сэмпла, поиск \(\alpha\) --- \(25\) итераций и \(4\) сэмпла, а validation prior predictive --- \(64\) сэмпла.

\subsection{Зачем нужен короткий общий старт}

Полностью случайный банк может в первые моменты давать почти неразличимые компоненты. Очень короткий общий прогрев создаёт разумные положительные интенсивности, но не должен успеть поглотить индивидуальную геометрию. После этого Pure и Wishart расходятся и обучаются независимо, что сохраняет честность сравнения.

\section{Предсказание после обучения}

\subsection{Совершенно новая траектория}

Для нового субъекта нет posterior по его индивидуальной матрице. Поэтому используется prior predictive:
\begin{equation}
p(s_{\mathrm{new}})=\sum_{k=1}^{K}\pi_k\int p(s_{\mathrm{new}}\mid k,U,\theta)\,p(U\mid k,\Omega_k)\,dU.
\label{eq:prior-predictive}
\end{equation}

На практике для каждого \(k\) сэмплируются матрицы \(U^{(r)}\sim p(U\mid k)\), вычисляются joint likelihood всей траектории и усредняются в log-domain через log-mean-exp. Одна матрица в одном Monte Carlo sample используется для всей траектории, иначе persistent random effect был бы разрушен.

\subsection{Уже наблюдавшаяся траектория}

Если доступен prefix конкретной траектории, по нему можно оценить \(q_{mk}(U)\) и \(\gamma_{mk}\), а затем предсказывать suffix:
\begin{equation}
p(s^{\mathrm{suffix}}\mid s^{\mathrm{prefix}})\approx\sum_k\gamma_{mk}(s^{\mathrm{prefix}})\int p(s^{\mathrm{suffix}}\mid U,k,s^{\mathrm{prefix}})q_{mk}(U\mid s^{\mathrm{prefix}})\,dU.
\end{equation}

Это posterior predictive. Он персонализирован: два субъекта одного кластера могут получить разные прогнозы благодаря разным posterior по \(U_m\).

\begin{warningbox}
Validation checkpoint выбирается по prior-predictive NLL, без локальной подгонки к validation или test траекториям. Posterior inference на test допустим только для явно обозначенных posterior-predictive и clustering diagnostics, но не должен незаметно попадать в метрику новой траектории.
\end{warningbox}

\section{Окна, общий parent и память}

\subsection{Правильное разбиение длинной траектории}

Если одну длинную траекторию разбить на окна \(X_{m1},\ldots,X_{mJ}\), эти окна не становятся независимыми субъектами. У них остаются общие \(z_m\) и \(U_m\):
\begin{equation}
p(X_{m1},\ldots,X_{mJ}\mid U_m,z_m)=\prod_{j=1}^{J}p(X_{mj}\mid U_m,z_m).
\label{eq:grouped-likelihood}
\end{equation}

Likelihood суммируется по окнам, но Wishart KL платится один раз за родителя, posterior \(q_m(U)\) общий, а COTIC state может сбрасываться на границе каждого окна. Так общий объём событий сохраняется, но доступная нейросети память контролируется экспериментально.

Если же завести отдельный \(U_{mj}\) для каждого окна, изменится сама модель: получится новый individual effect на каждой границе, хотя в DGP он постоянен.

\subsection{Невидимый prefix}

Для проверки переноса информации фиксируют один и тот же scored suffix, но меняют число видимых событий до его начала. Тогда можно сравнить четыре режима:
\begin{center}
\begin{tabular}{@{}lcc@{}}
\toprule
Режим & Neural state видит prefix & \(q(U,z)\) видит prefix \\
\midrule
Cold start & нет & нет, используется prior \\
Neural memory & да & нет \\
Latent memory & нет или reset & да \\
Full information & да & да \\
\bottomrule
\end{tabular}
\end{center}

Если latent-memory близок к full-information, persistent информация действительно переносится через \(U_m\). Если neural-memory не хуже, backbone успешно хранит её сам. Этот тест оценивает функциональную роль latent variable, а не буквальное восстановление параметров генератора.

\section{Что в модели идентифицируется}

Данные непосредственно наблюдают не \(\Omega_k\), \(\alpha\) или \(h_{\theta,k}\) по отдельности, а композицию
\begin{equation}
(h_{\theta,k},\alpha,U_m)\longmapsto\lambda_m(t).
\end{equation}

Гибкий backbone может частично перенести средний эффект \(U\) в собственные головы; изменение \(\alpha\) может компенсироваться масштабом \(U^{1/2}-I\); разные параметризации иногда дают близкие распределения интенсивностей. Поэтому большая сырая ошибка \(\Omega\) сама по себе не означает, что модель функционально неверна.

Главные метрики должны измерять наблюдаемое поведение:
\begin{itemize}
\item prior-predictive NLL новой траектории;
\item prefix-to-suffix NLL и способность к персонализированному прогнозу;
\item ошибка индивидуальной интенсивности \(\lambda_m(t)\);
\item ошибка ожидаемого числа событий и долей меток;
\item ARI и purity кластеризации;
\item устойчивость при разбиении одной траектории на reset-state окна.
\end{itemize}

Recovery \(\Omega_k\), correlation внедиагональных элементов, \(\kappa\) и KL полезны как вторичные диагностики оптимизации и выбранной системы координат.

\begin{intuitionbox}
Если две разные внутренние настройки музыкального синтезатора дают почти одинаковый звук, слушатель не может определить настройку по записи. Для нашей задачи «звук» --- распределение событий и интенсивностей, а \(\Omega\), \(\alpha\) и часть параметров backbone --- внутренние ручки синтезатора.
\end{intuitionbox}

\section{Как честно сравнивать Pure и Wishart}

Сравнение должно различаться только наличием random-effect ветки. Для DAN all-12 используется один split \(80/10/10\) с seed \(42\), одинаковые начальные веса backbone, одинаковый neural budget, одинаковая схема численного интегрирования и validation selection.

Pure Mixture и Wishart запускаются отдельными experiment runner. Отдельный Pure EM не нужен: если требуется endpoint-абляция внутри того же расписания, Wishart можно запустить с фиксированным \(\alpha=0\).

Нужно различать две NLL-метрики. Pure вычисляет обычный mixture likelihood. Wishart prior predictive интегрирует \(U\) по population prior. Нельзя сравнивать posterior-подогнанный Wishart с prior-only Pure и называть обе метрики одинаково.

\section{Типичные ошибки реализации}

\begin{enumerate}
\item \term{Новый \(U\) для каждого события или окна.} Это уничтожает persistent trajectory-level effect. Один Monte Carlo sample матрицы должен обслуживать всю оцениваемую траекторию или все окна одного parent.
\item \term{Независимый кластер для каждого окна.} Для grouped sequence переменная \(z_m\) общая, поэтому likelihood имеет вид \(\sum_k\pi_k\prod_jp(X_{mj}\mid k)\), а не произведение независимых смесей.
\item \term{Повторный posterior-множитель.} Факторизация \(p(s,z\mid U)=p(s\mid U)p(z\mid s,U)\) эквивалентна joint distribution; posterior нельзя дополнительно умножать в objective поверх \(\pi_kp(s\mid k,U)\).
\item \term{Неверное обновление \(\Omega\).} Формула \(CS/\tr S\) не является точным constrained M-step для objective~\eqref{eq:omega-mstep}.
\item \term{Обучение всех блоков одним Adam без разделения.} Это скрывает, кто компенсирует чью ошибку, и затрудняет диагностику. Базовый алгоритм явно разделяет local inference, population M-step, alpha search и neural M-step.
\item \term{Шум при validation.} Разные узлы Монте-Карло могут менять порядок близких контрольных точек. Для validation и test фиксируют один общий набор случайных узлов.
\item \term{Выбор checkpoint по test.} Все решения о cycle, \(\alpha\) и варианте модели принимаются на validation; test используется один раз для итоговой оценки.
\item \term{Потеря абсолютного времени при нарезке.} Если DGP или backbone использует фазу абсолютного времени, timestamps окон нельзя автоматически сдвигать к нулю без явного контроля.
\end{enumerate}

\section{Связь математики с объектами кода}

Архитектура реализации следует принципу «один класс --- одна задача». Основные соответствия приведены в таблице.

\begin{longtable}{@{}p{0.34\textwidth}p{0.58\textwidth}@{}}
\toprule
Класс или файл & Ответственность \\
\midrule
\endfirsthead
\toprule
Класс или файл & Ответственность \\
\midrule
\endhead
\code{TPPTrace} & Контракт между произвольным backbone и алгоритмом: event log-rates, quadrature или Monte Carlo nodes, weights и служебные маски. \\
\code{ActiveBlockDecoder} & Реализует формулу~\eqref{eq:active-rate}, component scores и prior-predictive Monte Carlo. \\
\code{VariationalParameters} & Хранит Cholesky-параметры \(Q_{mk}\) и параметр \(\kappa_{mk}\), гарантируя допустимость Wishart posterior. \\
\code{LocalWishartInference} & Оптимизирует free energy~\eqref{eq:free-energy} для локальных posterior. \\
\code{ResponsibilityUpdater} & Вычисляет softmax~\eqref{eq:responsibilities} и начальное balance-преобразование. \\
\code{PopulationMstep} & Решает spectral constrained-задачу~\eqref{eq:omega-mstep}. \\
\code{AlphaMstep} & Выполняет bounded scalar search по objective~\eqref{eq:alpha-objective}. \\
\code{PureMixtureTrainer} & Обучает baseline без random effect. \\
\code{ActiveTrainer} & Оркестрирует один цикл variational EM и выбирает лучший validation checkpoint. \\
\code{PureExperimentRunner} & Готовит данные, запускает и сохраняет отдельный Pure-эксперимент. \\
\code{WishartExperimentRunner} & Готовит данные, запускает и сохраняет отдельный Wishart-эксперимент. \\
\code{ExperimentRunner} & Абстрактная граница общего experiment workflow. \\
\bottomrule
\end{longtable}

Backbone-адаптеры находятся отдельно от Wishart inference. Это позволяет менять COTIC на THP, NHP или RMTPP, не переписывая математику random effect. Единственное требование --- получить корректный положительный банк базовых интенсивностей и trace для TPP likelihood.

\section{Формульная карта алгоритма}

\begin{summarybox}
\textbf{Базовый банк:}
\begin{equation*}
h_{\theta,k}(t\mid\History_t)\in\R_+^C.
\end{equation*}
\textbf{Скрытые переменные:}
\begin{equation*}
z_m\sim\Cat(\pi),\qquad U_m\mid z_m=k\sim\Wishart_C(\nu,\Omega_k/\nu),\qquad \tr\Omega_k=C.
\end{equation*}
\textbf{Decoder:}
\begin{equation*}
\lambda_m(t\mid k,U)=(1-\alpha)h_{\theta,k}(t)+\alpha\left[U^{1/2}\sqrt{h_{\theta,k}(t)}\right]^{\circ2}+10^{-6}.
\end{equation*}
\textbf{TPP log-likelihood:}
\begin{equation*}
\ell_{mk}(U)=\sum_i\log\lambda_{mk,c_i}(t_i;U)-\int_0^{T_m}\sum_c\lambda_{mkc}(t;U)dt.
\end{equation*}
\textbf{Variational posterior и free energy:}
\begin{equation*}
q_{mk}(U)=\Wishart_C(\kappa_{mk},Q_{mk}/\kappa_{mk}),\qquad \mathcal F_{mk}=-\E_q\ell_{mk}+\KL(q_{mk}\|p_k).
\end{equation*}
\textbf{Responsibilities:}
\begin{equation*}
\gamma_{mk}=\softmax_k(\log\pi_k-\mathcal F_{mk}).
\end{equation*}
\textbf{Population statistic и M-step:}
\begin{equation*}
S_k=\frac{\sum_m\gamma_{mk}Q_{mk}}{\sum_m\gamma_{mk}},\qquad \Omega_k=\arg\min_{\Omega\succ0,\tr\Omega=C}\bigl[\log|\Omega|+\tr(\Omega^{-1}S_k)\bigr].
\end{equation*}
\textbf{Pure endpoint:} при \(\alpha=0\) выполняется \(\lambda_m=h_{\theta,k}+10^{-6}\), и Wishart-ветка точно отключена.
\end{summarybox}

\section{Проверка понимания}

\begin{enumerate}
\item Чем интенсивность отличается от вероятности? \textit{Ответ: интенсивность имеет размерность «события/время»; вероятность на коротком интервале приблизительно равна интенсивности, умноженной на длину интервала.}
\item Почему в likelihood есть отрицательный интеграл? \textit{Ответ: он штрафует модель за обещанные события на всём интервале и является log-вероятностью отсутствия лишних событий.}
\item В чём различие между \(z_m\), \(U_m\) и \(\Omega_k\)? \textit{Ответ: \(z_m\) --- дискретный кластер, \(U_m\) --- индивидуальная матрица траектории, \(\Omega_k\) --- средняя геометрия population prior кластера.}
\item Почему матрица действует на \(\sqrt h\), а результат возводится в квадрат? \textit{Ответ: это даёт линейное смешивание амплитуд и одновременно гарантирует неотрицательность итоговой интенсивности.}
\item Что происходит при \(\alpha=0\)? \textit{Ответ: \(U\) исчезает из likelihood, и модель точно совпадает с Pure Mixture при том же постоянном floor.}
\item Зачем нужен KL в free energy? \textit{Ответ: он не позволяет локальному posterior произвольно подогнать каждую траекторию и связывает individual effects общей population law.}
\item Почему окна одного parent используют общий \(U\)? \textit{Ответ: \(U\) представляет постоянную характеристику субъекта, а граница вычислительного окна не создаёт нового субъекта.}
\item Почему raw \(\Omega\)-error вторична? \textit{Ответ: наблюдаемые данные определяют композицию backbone, \(\alpha\) и \(U\), поэтому внутренние параметры могут быть неединственными при близких функциональных интенсивностях.}
\end{enumerate}

\section{Заключение}

Active Block Wishart TPP --- это обучаемая вероятностная обёртка над произвольным нейросетевым временным точечным процессом. Backbone моделирует общие динамические закономерности каждого кластера, категориальная переменная \(z_m\) задаёт качественный тип траектории, а Wishart-переменная \(U_m\) добавляет устойчивую индивидуальную неоднородность внутри этого типа.

Главная конструктивная особенность модели --- точный endpoint \(\alpha=0\). Благодаря ему алгоритм может вести себя как чистая смесь, когда матричный random effect не поддерживается данными, и расширять пространство интенсивностей до стохастического семейства, когда persistent heterogeneity улучшает prediction.

Алгоритм обучения сохраняет это разделение ответственности: локальный E-шаг извлекает индивидуальные эффекты, exact population M-step обновляет \(\Omega_k\), одномерный поиск выбирает силу смешивания, а neural M-step адаптирует базовый банк. В результате каждую часть можно отдельно диагностировать, а качество модели следует оценивать прежде всего в наблюдаемом пространстве: по NLL, прогнозу suffix, интенсивностям, counts, marks и кластеризации.

\begin{thebibliography}{9}
\bibitem{daley} D. J. Daley, D. Vere-Jones. \textit{An Introduction to the Theory of Point Processes}. Springer, 2003.
\bibitem{rasmussen} J. G. Rasmussen. Temporal point processes: the conditional intensity function. \textit{Lecture Notes}, 2018.
\bibitem{muirhead} R. J. Muirhead. \textit{Aspects of Multivariate Statistical Theory}. Wiley, 1982.
\bibitem{blei} D. M. Blei, A. Kucukelbir, J. D. McAuliffe. Variational inference: a review for statisticians. \textit{Journal of the American Statistical Association}, 2017.
\bibitem{higham} N. J. Higham. \textit{Functions of Matrices: Theory and Computation}. SIAM, 2008.
\bibitem{cotic} V. Zholus et al. COTIC: continuous-time convolutional networks for event sequences. \textit{arXiv preprint}, 2022.
\end{thebibliography}
\end{document}
